% ============================================================
% CAPITOLO 21 - RIEPILOGO DELLA PARTE V
% ============================================================
\chapter{Riepilogo della Parte V}
\label{chap:riepilogo_v}

\section{Quadro delle note estranee}
\label{sec:sintesi_note_estranee}

\begin{center}
\small
\begin{tabularx}{\linewidth}{lLLLLL}
\toprule
\textbf{Nome} & \textbf{Tempo metrico} & \textbf{Raggiunta per...} & \textbf{Abbandonata per...} & \textbf{Preparata?} & \textbf{Dettaglio} \\
\midrule
Passaggio & debole & grado congiunto & grado congiunto & no & Sez. \ref{sec:nota_passaggio} \\
Volta & debole & grado congiunto & ritorno & no & Sez. \ref{sec:nota_volta} \\
Anticipazione & debole & qualsiasi & nota reale & no (è già risoluzione) & Sez. \ref{sec:anticipazione} \\
Appoggiatura & forte & salto (talvolta grado) & grado congiunto & no & Sez. \ref{sec:appoggiatura} \\
Sfuggita & debole & grado congiunto & salto opposto & no & Sez. \ref{sec:sfuggita} \\
Pedale & entrambi & tenuta & tenuta & sì & Sez. \ref{sec:pedale} \\
\bottomrule
\end{tabularx}
\end{center}

\section{Estranea o ritardo?}
\label{sec:estranea_o_ritardo}

Per distinguere una nota estranea da un ritardo conviene farsi tre domande.

\begin{enumerate}
	\item \textbf{La nota era già presente nell'accordo precedente, nella stessa voce, e viene tenuta?} Se sì può essere un ritardo (Sezione \ref{sec:definizione_ritardo}); se è raggiunta per salto o per grado senza preparazione è una nota estranea.
	\item \textbf{Cade sul tempo forte?} Il ritardo vi cade sempre. L'appoggiatura vi cade ma non è preparata; le altre note estranee cadono di solito sul debole.
	\item \textbf{Risolve per grado discendente?} Il ritardo risolve sempre in discesa; le note estranee possono anche salire o tornare indietro.
\end{enumerate}

Il ciclo generale della dissonanza (Sezione \ref{sec:ciclo_dissonanza}) vale in entrambi i casi.
